\documentclass[12pt]{article}
\usepackage[utf8]{inputenc}
\usepackage[T1]{fontenc}
\usepackage{hyperref}
\title{Venezuelan Sign Language Language Resources for K-12 Curriculum Development}
\author{Compiled by Claude Sonnet 5 (claude-sonnet-5) for Andr\'as Kornai}
\date{}
\begin{document}
\maketitle

\section{Introduction}
Venezuelan Sign Language (Lengua de Se\~nas Venezolana, LSV; ISO 639-3: vsl) is the language of the deaf community of Venezuela. This search found mainly academic work on deaf education. Wikipedia mentions a large dictionary published by the Federaci\'on Venezolana de Sordos and a national bilingual programme for LSV and Spanish. The search found no link to the dictionary, so it is not listed as a resource here.

\section{Primary Educational Resources}
\begin{itemize}
\item Article by Ana Mar\'ia Morales (UPEL-Instituto Pedag\'ogico de Caracas, 2008) that develops LSV study programmes for grades 1--3, with instructional games: \url{https://ve.scielo.org/pdf/edu/v12n41/art05.pdf}
\item UCV thesis ``Deditos M\'agicos'': audiovisual shorts for teaching LSV to deaf and hearing children: \url{https://saber.ucv.ve/jspui/handle/10872/1032}
\end{itemize}

\section{Linguistic and Reference Materials}
\begin{itemize}
\item Sapiens article on the norm in Venezuelan Sign Language and bilingual-bicultural planning: \url{https://biblat.unam.mx/en/revista/sapiens/articulo/la-norma-en-la-lengua-de-senas-venezolana}
\item Study of classroom interaction mediated by an LSV interpreter (P\'erez and Tovar, 2006): \url{https://dialnet.unirioja.es/descarga/articulo/2310278.pdf}
\item Further Dialnet article on deaf education: \url{https://dialnet.unirioja.es/descarga/articulo/7338762.pdf}
\item Wikipedia entry on the language: \url{https://wikipedia.com/wiki/ISO_639:vsl}
\end{itemize}

\section{Religious and Community Resources}
No religious resources were found. The Federaci\'on Venezolana de Sordos is named in the Wikipedia entry as publisher of a dictionary, but no website was seen in the results.

\section{Digital and Audio Resources}
The only digital resource found is the ``Deditos M\'agicos'' thesis above. No online dictionary or video corpus was located.

\section{Curriculum Development Context}
LSV has a research literature and some early-grade programme designs. Published classroom materials and a verifiable dictionary remain to be found, preferably through the Federaci\'on Venezolana de Sordos.
\end{document}